% ==============================================================================
% Capítulo 6: Conclusão
% ==============================================================================
\chapter{Conclusão}
\label{cap:conclusao}

\begin{instrucao}
\orientacao{O que vai na Conclusão}
A Conclusão fecha o trabalho. Ela retoma a questão de pesquisa e os objetivos,
mostra como cada um foi atingido, resume as contribuições, reconhece as
limitações e aponta trabalhos futuros. Como referência, de duas a quatro
páginas costumam ser suficientes.

Não apresente resultados novos nem repita a discussão do capítulo anterior. A
Conclusão costuma ter poucas citações, porque trata do que o próprio trabalho
produziu. Uma forma eficaz de organizá-la é percorrer os objetivos
específicos na mesma ordem em que aparecem na Introdução, indicando onde cada
um foi tratado.
\end{instrucao}

Este trabalho propôs uma heurística de \acs{ils} para o \acs{mkp} com
perturbação adaptativa e avaliou seu desempenho em instâncias de referência da
literatura. Os resultados mostraram que, com 60 segundos de execução, a
heurística obtém soluções mais próximas das melhores conhecidas do que o
solucionador HiGHS nas instâncias com 250 e 500 itens, respondendo
afirmativamente à questão de pesquisa formulada na
\cref{sec:introducao:problema}.

Os objetivos específicos foram atingidos da seguinte forma. A revisão das
abordagens heurísticas e das instâncias de referência foi apresentada no
\cref{cap:relacionados}. A heurística proposta e o modelo de \acs{pli} foram
descritos e implementados conforme o \cref{cap:metodologia}. A comparação com
o solucionador e com o \acs{ag} da literatura, bem como a avaliação da
perturbação adaptativa, foram apresentadas no \cref{cap:experimentos}.

% ------------------------------------------------------------------------------
\section{Contribuições e limitações}
\label{sec:conclusao:contribuicoes}

A principal contribuição do trabalho é uma heurística simples, com apenas dois
parâmetros e sem resolver programas lineares durante a busca, que obtém
resultados competitivos em tempo limitado. A análise dos componentes mostrou
que a perturbação adaptativa é responsável por boa parte do ganho de
desempenho, sobretudo nas instâncias com mais restrições.

Como limitações, a heurística ainda fica atrás dos melhores métodos híbridos
da literatura quando o tempo de execução não é restrito. Além disso, a
comparação com o \acs{ag} usou resultados publicados, obtidos em outro
hardware, e a avaliação se restringiu às instâncias da OR-Library.

% ------------------------------------------------------------------------------
\section{Trabalhos futuros}
\label{sec:conclusao:futuros}

\begin{instrucao}
\orientacao{Trabalhos futuros}
Evite listas genéricas como ``melhorar o desempenho'' ou ``testar com mais
dados''. Para cada sugestão, explique por que ela é promissora e como poderia
ser realizada, com base no que o trabalho mostrou.
\end{instrucao}

Os resultados sugerem três direções de continuidade. A primeira é usar a
relaxação linear apenas na construção da solução inicial, o que poderia
aproximar a heurística dos métodos híbridos sem aumentar o custo de cada
iteração. A segunda é adaptar a heurística a variantes do problema com
restrições adicionais, como conflitos entre itens, aproveitando o fato de que
a busca local mantém a viabilidade a cada movimento. A terceira é avaliar a
heurística em instâncias maiores e mais recentes, para verificar se a
vantagem sobre o solucionador se mantém com o crescimento do número de itens.
